% !TeX root = ../动手学强化学习笔记.tex
% 由主文档通过 \input 载入；请编译主文档。
\setcounter{chapter}{7} % 保留原书章号。
\chapter{时序差分学习}
\label{ch:rl-td}

时序差分学习（temporal-difference learning，TD）利用刚观察到的一步结果，同时借用当前后继价值。用已有估计构造学习目标称为自举（bootstrapping）。它使更新无需等待回合结束，却也让目标随估计一起变化。

\section{TD 误差与一步预测}

固定策略 $\pi$，令 $V_t$ 为当前价值表。一步 TD 的误差与更新为
\begin{equation}
\delta_t=R_{t+1}+\gamma V_t(S_{t+1})-V_t(S_t),\qquad
V_{t+1}(S_t)=V_t(S_t)+\alpha_t\delta_t.
\label{eq:rl-td-zero}
\end{equation}
其余表项保持不变。真实终止后取后继价值为零；采样截断不自动置零。给定采样前历史、当前状态和当前价值表，平均目标为 $(T_\pi V_t)(S_t)$，故平均更新方向是贝尔曼残差。除非后继估计已恰当，否则一步目标一般不是 $v_\pi(S_t)$ 的无偏样本。

例如当前估计为 $V(s)=3,V(s')=4$，观察奖励 $1$，取 $\gamma=0.9,\alpha=0.1$，则目标为 $4.6$、误差为 $1.6$，新估计为 $3.16$。当前更新承认“后继可能值得 $4$”，而不必真的等待该后继的全部未来奖励。

在有限折扣表格模型中，固定策略、奖励有界、有限初值、所有待估状态无限访问，且每个状态的步长由过去信息决定、位于 $(0,1]$ 并按其访问次数满足随机逼近级数条件时，标准 TD(0) 收敛定理给出 $V_t\to v_\pi$ 几乎必然成立\cite{sutton-barto}。该结论的严格证明属于随机逼近理论，本书不把期望更新的压缩性当作随机迭代完整证明。持续任务需确保所述访问条件确实成立，回合任务则需适当重置。

\section{Sarsa 与同策略控制}

Sarsa 名称对应状态、动作、奖励、下一状态、下一动作五个对象。先按当前行为策略选 $A_t$，观察转移，再按同一行为规则选 $A_{t+1}$，随后更新
\[
Q_{t+1}(S_t,A_t)=Q_t(S_t,A_t)+\alpha_t
\bigl[R_{t+1}+\gamma Q_t(S_{t+1},A_{t+1})-Q_t(S_t,A_t)\bigr].
\]
下一动作价值包含行为策略的探索影响；若后继是真实终止状态，则目标只剩奖励，不再采样下一动作。更新后执行已选出的 $A_{t+1}$，即可延续交互。

Sarsa 通常配合 $\varepsilon$-贪心策略实施广义策略迭代。固定探索率下，它学习带探索行为的后果；训练后直接改用贪心执行会改变所评价的策略。要得到趋向最优控制的标准结论，需要无限探索且极限趋于贪心等条件，有限训练中则通过独立评价检查实际策略。

\section{Q-learning 与期望 Sarsa}

Q 学习（Q-learning）使用最优备份目标
\[
Y_t^{\rm Q}=R_{t+1}+\gamma\max_aQ_t(S_{t+1},a),
\]
再按 $Q\leftarrow Q+\alpha(Y-Q)$ 更新访问的表项。行为策略可以选择非贪心动作，但目标假设下一步按当前最大值选择。期望 Sarsa（Expected Sarsa）则使用
\[
Y_t^{\rm ES}=R_{t+1}+\gamma\sum_a\pi_t(a\mid S_{t+1})Q_t(S_{t+1},a),
\]
显式平均下一动作随机性；其中 $\pi_t$ 是此次备份选定的目标策略，若与行为策略相同则为同策略版本。

\begin{example}[同一转移，三种目标]
当前 $Q(s,a)=3$，奖励为 $1$，后继两动作价值为 $4,2$，策略概率为 $1/4,3/4$；实际下一动作选到了价值为 $2$ 的动作。取 $\gamma=0.9$，Sarsa、Q-learning、期望 Sarsa 的目标分别为 $2.8,4.6,3.25$。若 $\alpha=0.1$，新表项分别为 $2.98,3.16,3.025$。差别来自目标策略和平均方式，而不是学习率。
\end{example}

有限状态动作、$0\leq\gamma<1$、奖励有界、有限初值、每个状态动作对无限访问，且逐对步长由过去信息决定、位于 $(0,1]$ 并满足随机逼近级数条件时，表格 Q-learning 几乎必然收敛到 $q_*$\cite{watkins-qlearning}。这既不保证任意探索规则满足访问条件，也不直接适用于神经网络和有限回放数据。

\section{最大化偏差与双重学习}

即使各动作估计无偏，通常也有 $\mathbb E[\max_a\widehat q(a)]\geq\max_a\mathbb E[\widehat q(a)]$。例如两动作真实价值均为零，独立估计各以相等概率取 $-1$ 或 $1$，最大值的期望为 $1/2$。选择最大值倾向挑中正噪声，这称为最大化偏差。

双重 Q 学习（Double Q-learning）维护 $Q^{(1)},Q^{(2)}$。若本次更新第一张表，先用它选择 $a^*\in\arg\max_aQ^{(1)}(S_{t+1},a)$，再用第二张表评价，目标为 $R_{t+1}+\gamma Q^{(2)}(S_{t+1},a^*)$；另一种更新对称进行。行为策略可对两表之和贪心并加探索。选择与评价分离缓解相同噪声被重复利用的问题，但有限样本和相关估计仍可产生其他偏差。

\section{习题与实验}
\begin{enumerate}
\item 将三目标例子中下一动作改为价值 $4$ 的动作，重新计算更新；说明哪两个算法的结果不变。
\item 在前章随机游走中比较 MC 与 TD(0) 的误差随交互步数的变化。明确价值误差对哪些状态平均，以及每条曲线有多少独立训练重复。
\item 在带悬崖的网格任务中比较 Sarsa 与 Q-learning，同时报告训练行为回报和独立贪心评价回报，解释两种指标回答的问题。
\end{enumerate}





